\Subhead{3.3.2}{Импульсные случайные процессы. Дробовой шум}

Если моменты событий $t_j$ образуют однородный пуассоновский поток интенсивности $\lambda$, то $\Pr(N_T=n)=\e^{-\lambda T}(\lambda T)^n/n!$, $\langle N_T\rangle=\operatorname{Var}N_T=\lambda T$. Линейно сформированный импульсный процесс
\begin{equation}
 i(t)=q\sum_jh(t-t_j),\qquad \int h(t)\dd t=1
\end{equation}
имеет $\overline i=q\lambda$. Независимость разных событий даёт ковариацию центрированного тока
\begin{equation}
 C_i(\tau)=q^2\lambda\int h(s+\tau)h(s)\dd s,
 \quad S_i(\omega)=q^2\lambda|H(\omega)|^2.
\end{equation}
При коротком импульсе $H\simeq1$ и $G_i(f)=2q\overline i$ — формула Шоттки в односторонней нормировке. Полный нецентрированный спектр содержит также линию постоянного тока. Случайная амплитуда импульсов заменяет $q^2$ на средний квадрат заряда, а не на квадрат среднего.

Пример — фототок идеального детектора когерентного света. Корреляции электронов или фотонов, мёртвое время и лавинное умножение нарушают простую пуассоновскую модель. Тепловой свет дополнительно даёт флуктуации интенсивности; его нельзя описать только независимым потоком одинаковых событий при любой временной шкале.
